\chapter[Kombinasyon · \textit{Temel}]{Kombinasyon}
\chaptermark{Kombinasyon}

Bölüm 4'te ele aldığımız permütasyon problemlerinde diziliş sırası belirleyiciydi: Bir podyumda birinci olmak ile üçüncü olmak arasındaki fark, bir kasanın şifresindeki rakamların dizilişi ya da bir yarışmadaki görev dağılımı doğrudan elemanların sırasına bağlıydı. 

Matematikte ve bilgisayar biliminde ise çoğu zaman nesnelerin hangi sırada dizildiğiyle değil, \textbf{yalnızca hangi nesnelerin seçildiğiyle} ilgileniriz:
\begin{itemize}
  \item $10$ kişilik bir sınıftan okul meclisine $3$ temsilci seçerken Ahmet, Ayşe ve Can'ın hangi sırayla belirlendiği temsilciler kurulunun yapısını değiştirmez. Kurul $\{\text{Ahmet}, \text{Ayşe}, \text{Can}\}$ kümesidir.
  \item Masadan çekilen $5$ iskambil kartının çekiliş sırası elin değerini etkilemez; elde toplanan kartlar kümesi belirleyicidir.
  \item Bir bilgisayar ağında $n$ sunucu arasından birbirine doğrudan bağlanacak iki sunucu belirlemek, sırasız bir ikili seçmektir.
\end{itemize}

Nesnelerin seçilme sırasının hesaba katılmadığı bu tür seçimlere \textbf{kombinasyon} (sırasız seçim) denir. Burada kombinasyon kavramını permütasyondan türetecek, $\binom{n}{r}$ binom katsayısını, simetri ve Pascal bağıntısını, ardından özdeş nesnelerin dağıtımında başvurulan ayraç yöntemini adım adım inceleyeceğiz.

% ============================================================
% 5.1 BİNOM KATSAYISI
% ============================================================
\section{Binom Katsayısı}

Küçük bir örnekle başlayalım: Dört kişilik bir arkadaş grubundan ($A, B, C, D$) iki kişilik bir proje ekibi kurmak istiyoruz. Kaç farklı ekip oluşturabiliriz?

Önce Bölüm 4'te gördüğümüz sıralı seçim (permütasyon) yaklaşımını uygulayalım:
\[ P(4, 2) = 4 \times 3 = 12 \]
Bu $12$ sıralı ikiliyi açıkça yazalım:
\begin{align*}
(A, B), \;\; (A, C), \;\; (A, D), \;\; (B, A), \;\; (B, C), \;\; (B, D), \\
(C, A), \;\; (C, B), \;\; (C, D), \;\; (D, A), \;\; (D, B), \;\; (D, C)
\end{align*}

Bir proje ekibi için $(A, B)$ ile $(B, A)$ arasında fark yoktur; her iki sıralı ikili de aynı iki kişiden oluşan $\{A, B\}$ ekibini belirtir. Benzer biçimde $(A, C)$ ile $(C, A)$ aynı ekiptir. 

Her ekip, seçilen $2$ kişinin kendi arasındaki $2! = 2$ farklı diziliminden ötürü listede tam \textbf{$2$ defa} yer almıştır. Farklı ekip sayısını bulmak için sıralı seçimlerin sayısını $2!$'e bölmeliyiz:
\[ \frac{P(4, 2)}{2!} = \frac{12}{2} = 6 \]
Elde edilen $6$ sırasız ekip şunlardır:
\[ \{A, B\}, \;\; \{A, C\}, \;\; \{A, D\}, \;\; \{B, C\}, \;\; \{B, D\}, \;\; \{C, D\} \]

Eğer bu $4$ kişi arasından $3$ kişilik bir ekip seçseydik, seçilen herhangi $3$ kişi kendi arasında $3! = 6$ farklı sırada dizilebilirdi. Dolayısıyla her ekip $P(4, 3) = 24$ sıralı üçlü listesinde tam $6$ defa tekrarlanacaktı. Gerçek ekip sayısı $\frac{P(4, 3)}{3!} = \frac{24}{6} = 4$ olurdu.

Bu gözlem bizi kombinasyonun formülüne ulaştırır:

\begin{definition}[Kombinasyon ve Binom Katsayısı]
$n$ ve $r$ iki doğal sayı ve $0 \le r \le n$ olsun. $n$ elemanlı bir kümenin $r$ elemanlı her bir alt kümesine $n$'nin $r$'li \textbf{kombinasyonu} denir. Bu alt kümelerin toplam sayısı $\binom{n}{r}$ (veya $C(n, r)$) simgesiyle gösterilir ve ``$n$'nin $r$'lisi'' diye okunur.
\end{definition}

\begin{theorem}
Her $0 \le r \le n$ için:
\[ \binom{n}{r} = \frac{P(n, r)}{r!} = \frac{n!}{r!(n - r)!} \]
\end{theorem}
\begin{proof}
$n$ elemanlı bir kümeden $r$ elemanlı sıralı bir dizi ($P(n, r)$) oluşturma sürecini iki aşamalı bir işlem olarak kurgulayalım:
\begin{enumerate}
  \item \textbf{1. Aşama (Kümeyi seçmek):} Sırasına bakılmaksızın hangi $r$ elemanın seçileceğini belirleriz. Bu seçim tanım gereği $\binom{n}{r}$ farklı şekilde yapılabilir.
  \item \textbf{2. Aşama (Seçilenleri dizmek):} Seçilen $r$ elemanı bir sıraya dizeriz. $r$ farklı nesne kendi arasında $r!$ farklı şekilde sıralanır.
\end{enumerate}
Bölüm 3'teki çarpma kuralı uyarınca bu iki aşamanın çarpımı toplam sıralı seçimlerin sayısını ($P(n, r)$) verir:
\[ \binom{n}{r} \times r! = P(n, r) \]
Eşitliğin her iki tarafını $r!$'e bölersek:
\[ \binom{n}{r} = \frac{P(n, r)}{r!} \]
elde edilir. Bölüm 4'teki $P(n, r) = \frac{n!}{(n-r)!}$ ifadesi yerine konduğunda:
\[ \binom{n}{r} = \frac{n!}{r!(n - r)!} \]
bulunur.
\end{proof}

Sınır durumları inceleyelim:
\begin{itemize}
  \item $\binom{n}{0} = \frac{n!}{0!(n-0)!} = \frac{n!}{1 \cdot n!} = 1$ ($n$ elemandan $0$ eleman seçmenin $1$ yolu vardır: boş küme $\emptyset$).
  \item $\binom{n}{1} = \frac{n!}{1!(n-1)!} = \frac{n \cdot (n-1)!}{(n-1)!} = n$ ($n$ nesneden tek bir nesne seçmenin $n$ yolu vardır).
  \item $\binom{n}{n} = \frac{n!}{n!(n-n)!} = \frac{n!}{n! \cdot 0!} = 1$ ($n$ nesnenin tamamını seçmenin tek bir yolu vardır: kümenin kendisi).
\end{itemize}

\subsection{Kümeler Kuramı ile Bağlantı}

Bölüm 1'de $n$ elemanlı bir kümenin toplam $2^n$ adet alt kümesi olduğunu çarpma kuralıyla doğrulamıştık. Kombinasyon bu alt kümeleri eleman sayılarına göre ayrık sınıflara böler:
\begin{itemize}
  \item $0$ elemanlı alt küme sayısı: $\binom{n}{0}$
  \item $1$ elemanlı alt küme sayısı: $\binom{n}{1}$
  \item $2$ elemanlı alt küme sayısı: $\binom{n}{2}$
  \item \dots
  \item $n$ elemanlı alt küme sayısı: $\binom{n}{n}$
\end{itemize}
Bu alt küme sınıfları birbirini dışlayan ayrık durumlar olduğundan, Bölüm 3'teki toplama kuralı gereğince tüm alt kümelerin toplamı:
\[ \sum_{r=0}^n \binom{n}{r} = \binom{n}{0} + \binom{n}{1} + \binom{n}{2} + \dots + \binom{n}{n} = 2^n \]
eşitliğini sağlar. Bu temel özdeşliği Bölüm 6'da binom teoremi çerçevesinde yeniden ele alacağız.

\subsection{Çözümlü Örnekler}

\begin{example}
Aşağıdaki kombinasyonların değerlerini hesaplayınız:
a) $\binom{7}{3}$ \quad b) $\binom{10}{2}$ \quad c) $\binom{8}{4}$
\end{example}
\begin{proof}[Çözüm]
Kombinasyon hesaplarken büyük faktöriyelleri açmak yerine, payı $r$ adım geriye açıp paydaya $r!$ yazmak daha pratiktir:
\begin{itemize}
  \item a) $\binom{7}{3} = \frac{7 \times 6 \times 5}{3 \times 2 \times 1} = \frac{210}{6} = 35$.
  \item b) $\binom{10}{2} = \frac{10 \times 9}{2 \times 1} = \frac{90}{2} = 45$.
  \item c) $\binom{8}{4} = \frac{8 \times 7 \times 6 \times 5}{4 \times 3 \times 2 \times 1} = \frac{8 \times 7 \times 6 \times 5}{24} = 7 \times 2 \times 5 = 70$.
\end{itemize}
\end{proof}

\begin{example}
$12$ üyeli bir öğrenci topluluğundan $4$ kişilik bir disiplin kurulu kaç farklı şekilde seçilebilir?
\end{example}
\begin{proof}[Çözüm]
Kurul üyeleri arasında unvan ya da görev farkı bulunmadığından seçim sırası önemsizdir:
\[ \binom{12}{4} = \frac{12 \times 11 \times 10 \times 9}{4 \times 3 \times 2 \times 1} \]
Paydadaki $4 \times 3 = 12$ paydaki $12$ ile sadeleşir. Kalan $2$ ise $10$'u sadeleştirip $5$ yapar:
\[ = 11 \times 5 \times 9 = 495 \]
farklı şekilde kurul seçilebilir.
\end{proof}

\begin{example}
Bir çember üzerinde birbirinden farklı $8$ nokta işaretlenmiştir.
\begin{enumerate}
  \item Köşeleri bu noktalardan seçilen kaç farklı doğru parçası (kiriş) çizilebilir?
  \item Köşeleri bu noktalardan seçilen kaç farklı üçgen oluşturulabilir?
\end{enumerate}
\end{example}
\begin{proof}[Çözüm]
Çember üzerindeki noktaların hiçbir üçü doğrusal değildir.
\begin{enumerate}
  \item Bir doğru parçası belirlemek için bu $8$ noktadan herhangi $2$ tanesini seçmek yeterlidir. Noktaların seçilme sırası doğru parçasını değiştirmez ($AB$ ile $BA$ aynı kiriştir):
  \[ \binom{8}{2} = \frac{8 \times 7}{2 \times 1} = 28 \text{ kiriş çizilebilir.} \]
  
  \item Bir üçgen belirlemek için bu $8$ noktadan herhangi $3$ tanesini seçmek gerekir:
  \[ \binom{8}{3} = \frac{8 \times 7 \times 6}{3 \times 2 \times 1} = 56 \text{ farklı üçgen çizilebilir.} \]
\end{enumerate}
\end{proof}

\begin{example}
Dışbükey (konveks) bir $n$ kenarlı çokgenin toplam kaç köşegeni vardır? Formülü kombinasyon yardımıyla elde ediniz.
\end{example}
\begin{proof}[Çözüm]
$n$ kenarlı bir çokgenin $n$ adet köşesi bulunur. Bu $n$ köşe arasından seçilecek herhangi iki köşe bir doğru parçası tanımlar:
\[ \text{Tüm Doğru Parçaları} = \binom{n}{2} \]
Bu doğru parçaları iki gruptan oluşur:
\begin{itemize}
  \item Çokgenin kenarları (komşu iki köşeyi birleştiren $n$ adet doğru parçası).
  \item Köşegenler (komşu olmayan iki köşeyi birleştiren doğru parçaları).
\end{itemize}
Dolayısıyla köşegen sayısı, tüm doğru parçalarının sayısından kenar sayısı çıkarılarak bulunur:
\[ \text{Köşegen Sayısı} = \binom{n}{2} - n = \frac{n(n - 1)}{2} - n = \frac{n^2 - n - 2n}{2} = \frac{n(n - 3)}{2} \]
Örneğin bir altıgen için ($n = 6$): $\frac{6 \times 3}{2} = 9$ köşegen elde edilir.
\end{proof}


% ============================================================
% 5.2 SİMETRİ VE PASCAL BAĞINTISI
% ============================================================
\section{Simetri ve Pascal Bağıntısı}

Kombinasyon katsayıları belirgin cebirsel ve simetrik özellikler taşır. Bu özelliklerin ilkini basit bir soruyla görebiliriz:

$10$ kişilik bir sınıftan kampa gidecek $8$ öğrenciyi mi seçmek daha kolaydır, yoksa geride kalacak $2$ öğrenciyi mi?

Kampa gidecek $8$ kişiyi belirlemek, geride kalacak $2$ kişiyi belirlemekle aynı seçim işlemidir. Gidecekler belirlendiğinde, kalanlar da kendiliğinden belirlenir. Formülle inceleyelim:
\[ \binom{10}{8} = \frac{10!}{8! \times 2!} = 45 \quad \text{ve} \quad \binom{10}{2} = \frac{10!}{2! \times 8!} = 45 \]
İki değer birbirine eşittir.

\begin{theorem}[Simetri Özelliği]
Her $0 \le r \le n$ için:
\[ \binom{n}{r} = \binom{n}{n - r} \]
\end{theorem}
\begin{proof}
Formülü doğrudan uygulayalım:
\[ \binom{n}{n - r} = \frac{n!}{(n - r)! \times (n - (n - r))!} = \frac{n!}{(n - r)! \times r!} \]
Çarpmanın değişme özelliği gereği paydadaki $(n - r)! \times r! = r! \times (n - r)!$'dir. Buradan:
\[ \binom{n}{n - r} = \frac{n!}{r!(n - r)!} = \binom{n}{r} \]
sonucuna varılır.
\end{proof}

Bu özellik hesaplamaları önemli ölçüde sadeleştirir. Örneğin $\binom{100}{98}$ ifadesi doğrudan simetri kuralıyla açılır:
\[ \binom{100}{98} = \binom{100}{100 - 98} = \binom{100}{2} = \frac{100 \times 99}{2} = 4950 \]

\subsection{Pascal Bağıntısı: Bir Elemanın Durumu}

Şimdi kombinatorik ispatların klasik yaklaşımlarından birini inceleyelim. $n$ kişilik bir gruptan $r$ kişilik bir komite kurmak istiyoruz ($\binom{n}{r}$). 

Bu gruptaki belirli bir kişiye odaklanalım; bu kişi Ahmet olsun. Olası tüm komiteleri Ahmet'in durumuna göre iki ayrık duruma ayırabiliriz:
\begin{enumerate}
  \item \textbf{1. Durum (Ahmet'in yer aldığı komiteler):} Ahmet komitededir. Komitenin kalan $r - 1$ üyesi, gruptaki diğer $n - 1$ kişi arasından seçilir:
  \[ \binom{n - 1}{r - 1} \text{ seçenek.} \]
  \item \textbf{2. Durum (Ahmet'in yer almadığı komiteler):} Ahmet komitede yer almaz. Dolayısıyla $r$ kişilik komitenin tamamı diğer $n - 1$ kişi arasından seçilmelidir:
  \[ \binom{n - 1}{r} \text{ seçenek.} \]
\end{enumerate}

Ahmet bir komitede ya vardır ya da yoktur; bu iki durum ayrıktır. Toplama kuralı gereğince iki durumun toplamı tüm komitelerin sayısını verir.

\begin{theorem}[Pascal Bağıntısı]
Her $1 \le r \le n$ için:
\[ \binom{n}{r} = \binom{n - 1}{r - 1} + \binom{n - 1}{r} \]
\end{theorem}
\begin{proof}
Yukarıdaki kombinatorik akıl yürütme doğrudan bir ispattır. Cebirsel olarak doğrulamak için sağ taraftaki terimleri ortak paydaya getirelim:
\[ \binom{n - 1}{r - 1} + \binom{n - 1}{r} = \frac{(n-1)!}{(r-1)!(n-r)!} + \frac{(n-1)!}{r!(n-1-r)!} \]
Birinci kesri $r$ ile, ikinci kesri $(n - r)$ ile genişletelim:
\begin{align*}
&= \frac{r \cdot (n-1)!}{r!(n-r)!} + \frac{(n-r) \cdot (n-1)!}{r!(n-r)!} \\
&= \frac{[r + (n - r)] \cdot (n-1)!}{r!(n-r)!} \\
&= \frac{n \cdot (n-1)!}{r!(n-r)!} = \frac{n!}{r!(n-r)!} = \binom{n}{r}
\end{align*}
böylece eşitlik gösterilmiş olur.
\end{proof}

\subsection{Çözümlü Örnekler}

\begin{example}
$\binom{10}{3} + \binom{10}{4} + \binom{11}{5}$ toplamını tek bir binom katsayısı olarak yazınız.
\end{example}
\begin{proof}[Çözüm]
Pascal bağıntısını adım adım uygulayalım.
İlk iki terimi ele alalım: $\binom{10}{3} + \binom{10}{4}$. 
Pascal kuralı gereğince $\binom{n-1}{r-1} + \binom{n-1}{r} = \binom{n}{r}$ olduğundan:
\[ \binom{10}{3} + \binom{10}{4} = \binom{11}{4} \]
olur. Bu sonucu üçüncü terimle birleştirelim:
\[ \binom{11}{4} + \binom{11}{5} \]
Tekrar Pascal kuralı uygulanırsa ($n = 12, r = 5$):
\[ \binom{11}{4} + \binom{11}{5} = \binom{12}{5} \]
elde edilir.
\end{proof}

\begin{example}[Kombinatorik İspat: Başkanlı Komite Modeli]
Her $1 \le r \le n$ için aşağıdaki eşitliğin doğruluğunu gösteriniz:
\[ n \binom{n - 1}{r - 1} = r \binom{n}{r} \]
\end{example}
\begin{proof}[Çözüm]
Bu eşitlik cebirsel yoldan doğrulanabileceği gibi, Bölüm 10'da ayrıntılandıracağımız iki yoldan sayma (çifte sayım) yöntemiyle de kurulabilir.

$n$ kişi arasından $r$ kişilik bir komite ve bu komiteye bir \textbf{başkan} seçeceğimizi varsayalım. Bu seçim iki farklı sırayla yapılabilir:

\begin{itemize}
  \item \textbf{1. Yol (Önce başkan, sonra üyeler):}
  $n$ aday arasından başkan seçilir ($n$ seçenek).
  Kalan $n - 1$ kişi arasından komitenin diğer $r - 1$ üyesi belirlenir ($\binom{n-1}{r-1}$ seçenek).
  Çarpma kuralıyla:
  \[ n \times \binom{n-1}{r-1} \]
  
  \item \textbf{2. Yol (Önce komite, sonra başkan):}
  $n$ kişi arasından $r$ kişilik komite oluşturulur ($\binom{n}{r}$ seçenek).
  Seçilen $r$ kişi arasından bir başkan belirlenir ($r$ seçenek).
  Çarpma kuralıyla:
  \[ \binom{n}{r} \times r \]
\end{itemize}
Her iki sayım da aynı kümenin eleman sayısını verdiğinden:
\[ n \binom{n - 1}{r - 1} = r \binom{n}{r} \]
eşitliği sağlanır.
\end{proof}

\begin{example}[C Dilinde Pascal Bağıntısıyla Rekürsif Kombinasyon]
Pascal bağıntısını kullanarak kombinasyon hesaplayan bir C fonksiyonu yazınız ve algoritmayı açıklayınız.
\end{example}
\begin{proof}[Çözüm]
Pascal bağıntısı $\binom{n}{r} = \binom{n-1}{r-1} + \binom{n-1}{r}$ doğal bir rekürsiyon (rekürsiyon) yapısı sunar. Taban durumları ise $r = 0$ ve $r = n$ olduğunda değerin $1$ olmasıdır:

\begin{lstlisting}[language=CTemel]
// Pascal bagintisi ile kombinasyon hesabi
long long kombinasyon(int n, int r) {
    // Taban durumlari (base cases)
    if (r == 0 || r == n) {
        return 1;
    }
    // Ozyinelemeli adim: Pascal kurali
    return kombinasyon(n - 1, r - 1) + kombinasyon(n - 1, r);
}
\end{lstlisting}

Bu fonksiyon matematiksel tanımı doğrudan uygular; ancak büyük $n$ değerlerinde aynı alt problemleri defalarca hesaplar. Bu verimsizlik Bölüm 26'da dinamik programlama yaklaşımıyla giderilecektir.
\end{proof}


% ============================================================
% 5.3 AYRAÇ YÖNTEMİ (STARS AND BARS)
% ============================================================
\section{Ayraç Yöntemi (Stars and Bars)}

Şimdiye kadar ele aldığımız kombinasyon problemlerinde üzerinde çalışılan nesneler birbirinden ayırt edilebilirdi. Dağıtılacak nesnelerin \textbf{birbiriyle tamamen özdeş} olduğu durumlarda sayım nasıl yapılır?

Somut bir problem ele alalım: Birbirinin tıpatıp aynısı olan $7$ elmayı $3$ çocuğa ($A, B, C$) kaç farklı şekilde paylaştırabiliriz?

Elmalar özdeş olduğundan hangi elmanın kime verildiği değil, \textbf{her çocuğun kaç elma aldığı} önemlidir.
$A$'nın aldığı elma sayısına $x_1$, $B$'ninkine $x_2$, $C$'ninkine $x_3$ dersek, problem şu denkleme denktir:
\[ x_1 + x_2 + x_3 = 7 \quad (x_1, x_2, x_3 \in \mathbb{N}, \;\; x_i \ge 0) \]

Bu denklemin çözümlerini tek tek listelemek yerine, kombinatorikte sıkça başvurulan \textbf{ayraç yöntemi} (\textit{Stars and Bars}) modelini kullanırız.

\subsection{Yöntemin Mantığı}

Yöntemin İngilizce adı (\textit{stars and bars}) görseli de açıklar: özdeş nesneleri yıldız ($\star$), grupları ayıran işaretleri ise düşey çizgi ($|$) ile temsil ederiz. $7$ adet özdeş elmayı yan yana dizilmiş $7$ adet yıldız ($\star$) ile gösterelim:
\[ \star \quad \star \quad \star \quad \star \quad \star \quad \star \quad \star \]
Bu $7$ yıldızı $3$ gruba ayırmak için araya \textbf{$3 - 1 = 2$ adet ayraç ($|$)} yerleştirmek yeterlidir.

Örnek dizilimler:
\begin{itemize}
  \item $\star \; \star \mid \star \; \star \; \star \mid \star \; \star \implies x_1 = 2, \;\; x_2 = 3, \;\; x_3 = 2$
  \item $\mid \star \; \star \; \star \; \star \mid \star \; \star \; \star \implies x_1 = 0, \;\; x_2 = 4, \;\; x_3 = 3$ (ilk çocuk hiç elma almadı)
  \item $\star \; \star \; \star \; \star \; \star \; \star \; \star \mid \mid \implies x_1 = 7, \;\; x_2 = 0, \;\; x_3 = 0$ (tüm elmaları ilk çocuk aldı)
\end{itemize}

Her paylaşım, $7$ yıldız ($\star$) ve $2$ ayraçtan ($|$) oluşan toplam $7 + 2 = 9$ karakterlik bir dizilimle birebir eşleşir.

Soru şuna indirgenir: Yan yana duran $9$ boş pozisyondan hangi $2$ tanesine ayraç ($|$) yerleştirilecektir?
Geriye kalan yerlere yıldızlar gelecektir. Sıra önemsiz olduğundan bu doğrudan bir kombinasyondur:
\[ \binom{9}{2} = \frac{9 \times 8}{2 \times 1} = 36 \]
Demek ki $7$ özdeş elma $3$ çocuğa $36$ farklı şekilde dağıtılabilir.

\begin{theorem}[Ayraç Yöntemi Teoremi — Negatif Olmayan Çözümler]
$n$ adet özdeş nesnenin $k$ farklı kutuya dağıtılma biçimlerinin (veya $x_1 + x_2 + \dots + x_k = n$ denkleminin $x_i \ge 0$ tam sayı çözümlerinin) toplam sayısı:
\[ \binom{n + k - 1}{k - 1} \]
formülüyle hesaplanır.
\end{theorem}
\begin{proof}
$n$ adet yıldız ve kutuları birbirinden ayırmak için $k - 1$ adet ayraç kullanılır. Toplam pozisyon sayısı $n + (k - 1) = n + k - 1$'dir. Bu pozisyonlardan ayraçların yerleşeceği $k - 1$ konumu belirlemenin yolu $\binom{n+k-1}{k-1}$ tanedir.
\end{proof}

\subsection{Her Kutuya En Az Bir Nesne Düşmesi Durumu (\texorpdfstring{$x_i \ge 1$}{x\_i >= 1})}

Her çocuğun \textbf{en az bir elma} alması koşulu konursa, yani $x_1 + x_2 + \dots + x_k = n$ denkleminin \textbf{pozitif tam sayı} ($x_i \ge 1$) çözümleri aranırsa iki yoldan ilerlenebilir:

\textbf{1. Yol (Önden Dağıtma):}
$n$ elmadan önce her çocuğa birer tane verilir. Toplam $k$ elma dağıtılmış olur ve geriye $n - k$ elma kalır. Kalan $n - k$ elma çocuklara serbestçe (negatif olmayan biçimde) paylaştırılabilir:
\[ \binom{(n - k) + k - 1}{k - 1} = \binom{n - 1}{k - 1} \]

\textbf{2. Yol (Aralıklara Ayraç Koyma):}
$n$ adet yıldız yan yana dizildiğinde aralarında tam $n - 1$ adet boşluk bulunur:
\[ \star \;\; \underbrace{\quad}_{\text{boşluk}} \;\; \star \;\; \underbrace{\quad}_{\text{boşluk}} \;\; \star \;\; \dots \;\; \star \]
Hiçbir kutunun boş kalmaması, iki ayracın aynı boşluğa konmaması ve uçlara ayraç yerleştirilmemesi demektir. Dolayısıyla $k - 1$ adet ayraç, bu $n - 1$ boşluk arasından seçilen yerlere konur:
\[ \binom{n - 1}{k - 1} \]

\begin{theorem}[Pozitif Tam Sayı Çözümleri]
$x_1 + x_2 + \dots + x_k = n$ denkleminin pozitif tam sayılardaki ($x_i \ge 1$) çözüm sayısı:
\[ \binom{n - 1}{k - 1} \]
formülüyle verilir.
\end{theorem}

\subsection{Çözümlü Örnekler}

\begin{example}
Bir pastanede $4$ çeşit kurabiye (çikolatalı, fındıklı, tarçınlı, vanilyalı) satılmaktadır ve her çeşitten yeterince vardır. Bir müşteri toplam $10$ adet kurabiye seçerek bir paket yaptırmak istediğinde kaç farklı paket oluşturabilir?
\end{example}
\begin{proof}[Çözüm]
Kurabiyeler aynı çeşit içinde özdeştir. Seçilecek miktarlara $x_1, x_2, x_3, x_4$ dersek:
\[ x_1 + x_2 + x_3 + x_4 = 10 \quad (x_i \ge 0) \]
Burada $n = 10$ nesne ve $k = 4$ farklı kategori söz konusudur.
Ayraç yöntemi teoremini uygulayalım:
\[ \binom{n + k - 1}{k - 1} = \binom{10 + 4 - 1}{4 - 1} = \binom{13}{3} \]
Değeri hesaplayalım:
\[ \binom{13}{3} = \frac{13 \times 12 \times 11}{3 \times 2 \times 1} = \frac{13 \times 12 \times 11}{6} = 13 \times 2 \times 11 = 286 \]
Farklı paket seçeneklerinin sayısı $286$'dır.
\end{proof}

\begin{example}[Değişken Değiştirme Tekniği]
$a + b + c = 18$ denkleminin $a \ge 2$, $b \ge 3$ ve $c \ge 0$ koşullarını sağlayan kaç farklı tam sayı çözümü vardır?
\end{example}
\begin{proof}[Çözüm]
Kısıtları standart $\ge 0$ biçimine dönüştürmek için değişken değiştirme uygularız:
\begin{itemize}
  \item $a \ge 2 \implies a' = a - 2 \ge 0 \implies a = a' + 2$
  \item $b \ge 3 \implies b' = b - 3 \ge 0 \implies b = b' + 3$
  \item $c \ge 0 \implies c' = c \ge 0 \implies c = c'$
\end{itemize}
Bu tanımları denklemde yerine koyalım:
\[ (a' + 2) + (b' + 3) + c' = 18 \implies a' + b' + c' + 5 = 18 \implies a' + b' + c' = 13 \]
Böylece $n = 13$ ve $k = 3$ parametrelerine sahip standart biçime ulaşırız:
\[ \binom{13 + 3 - 1}{3 - 1} = \binom{15}{2} = \frac{15 \times 14}{2} = 105 \]
farklı çözüm vardır.
\end{proof}

\begin{example}[Üstten Kısıtlı Dağıtım]
$x_1 + x_2 + x_3 = 9$ denkleminin $0 \le x_i \le 4$ kısıtını sağlayan tam sayı çözüm sayısını bulunuz.
\end{example}
\begin{proof}[Çözüm]
Değişkenlerin alabileceği en büyük değer $4$ olarak sınırlandırılmıştır. Çözümü tümleyen küme üzerinden kuralım:
\[ \text{İstenen} = \text{Tüm Çözümler} - \text{En az bir değişkenin } \ge 5 \text{ olduğu durumlar} \]

\begin{enumerate}
  \item \textbf{Tüm Çözümler ($x_i \ge 0$):}
  \[ \binom{9 + 3 - 1}{3 - 1} = \binom{11}{2} = \frac{11 \times 10}{2} = 55 \]
  
  \item \textbf{Bir değişkenin $\ge 5$ olması:}
  Değişkenlerden biri (örneğin $x_1$) en az $5$ olsun: $x_1 \ge 5$.
  $x_1' = x_1 - 5 \ge 0$ dönüşümüyle:
  \[ x_1' + x_2 + x_3 = 9 - 5 = 4 \]
  Bu denklemin çözüm sayısı:
  \[ \binom{4 + 3 - 1}{3 - 1} = \binom{6}{2} = 15 \]
  En az $5$ olacak değişkeni seçmenin $\binom{3}{1} = 3$ yolu vardır ($x_1, x_2$ veya $x_3$).
  
  \item \textbf{İki değişken birden $\ge 5$ olabilir mi?}
  İki değişken aynı anda $\ge 5$ olsaydı toplam en az $5 + 5 = 10$ olurdu; oysa denklem toplamı $9$'dur. Dolayısıyla bu durum gerçekleşemez.
\end{enumerate}

Koşulu bozan durumların sayısı $3 \times 15 = 45$'tir.
O hâlde aranan çözüm sayısı:
\[ 55 - 45 = 10 \]
olur.
\end{proof}


% ============================================================
% 5.4 KARIŞIK SEÇME PROBLEMLERİ
% ============================================================
\section{Karışık Seçme Problemleri}

Kombinatorik problemlerinde çoğu zaman kombinasyon, permütasyon ve durum analizi bir arada kullanılır. Temel yaklaşım kuralı şudur: \textbf{Önce seçim yapılır (kombinasyon), gerek duyuluyorsa ardından sıralama uygulanır (permütasyon).}

\subsection{Çözümlü Örnekler}

\begin{example}
$6$ erkek ve $5$ kadın arasından $4$ kişilik bir temsilci heyeti seçilecektir.
\begin{enumerate}
  \item Heyette \textbf{en az $2$ kadın} bulunması şartıyla kaç farklı seçim yapılabilir?
  \item Belirli iki kişinin (Ahmet ve Zeynep) \textbf{birlikte bulunmaması} şartıyla kaç farklı heyet kurulabilir?
\end{enumerate}
\end{example}
\begin{proof}[Çözüm]
\begin{enumerate}
  \item Heyetteki kadın sayısı $2, 3$ veya $4$ olabilir. Bu durumlar ayrık olduğundan ayrı ayrı hesaplanıp toplanır:
  \begin{itemize}
    \item \textbf{2 Kadın, 2 Erkek:} $\binom{5}{2} \times \binom{6}{2} = 10 \times 15 = 150$
    \item \textbf{3 Kadın, 1 Erkek:} $\binom{5}{3} \times \binom{6}{1} = 10 \times 6 = 60$
    \item \textbf{4 Kadın, 0 Erkek:} $\binom{5}{4} \times \binom{6}{0} = 5 \times 1 = 5$
  \end{itemize}
  Toplama kuralı gereğince:
  \[ 150 + 60 + 5 = 215 \text{ farklı heyet kurulabilir.} \]

  \item Tümleyen yaklaşımını kullanalım:
  \[ \text{İstenen} = \text{Tüm Heyetler} - \text{Ahmet ile Zeynep'in Bulunduğu Heyetler} \]
  \begin{itemize}
    \item Toplam $6 + 5 = 11$ kişi vardır. Kısıt olmaksızın $4$ kişi:
    \[ \binom{11}{4} = \frac{11 \times 10 \times 9 \times 8}{4 \times 3 \times 2 \times 1} = 330 \]
    \item Ahmet ve Zeynep'in ikisinin de seçildiği durumlar: Bu iki kişi ayrıldıktan sonra kalan $11 - 2 = 9$ kişiden komiteyi tamamlayacak $4 - 2 = 2$ kişi seçilir:
    \[ \binom{9}{2} = \frac{9 \times 8}{2} = 36 \]
  \end{itemize}
  Buna göre Ahmet ve Zeynep'in birlikte yer almadığı heyet sayısı:
  \[ 330 - 36 = 294 \]
  olur.
\end{enumerate}
\end{proof}

\begin{example}[Izgarada Dikdörtgen Sayma]
Yatayda $5$ paralel doğru ile düşeyde $7$ paralel doğru birbirini dik kesmektedir. Kenarları bu doğrular üzerinde bulunan kaç farklı dikdörtgen oluşur?
\end{example}
\begin{proof}[Çözüm]
Bir dikdörtgen tanımlamak için:
\begin{itemize}
  \item Alt ve üst kenarı belirleyecek \textbf{$2$ adet yatay doğru},
  \item Sol ve sağ kenarı belirleyecek \textbf{$2$ adet düşey doğru}
\end{itemize}
seçmek gerekir. Seçilen her yatay doğru çifti ile düşey doğru çifti tek bir dikdörtgen sınırlar:
\begin{itemize}
  \item $5$ yatay doğrudan $2$ tanesini seçmenin yolu: $\binom{5}{2} = 10$.
  \item $7$ düşey doğrudan $2$ tanesini seçmenin yolu: $\binom{7}{2} = 21$.
\end{itemize}
Çarpma kuralı uyarınca toplam dikdörtgen sayısı:
\[ \binom{5}{2} \times \binom{7}{2} = 10 \times 21 = 210 \]
olur. Genel olarak $m$ yatay ve $n$ düşey doğru arasında $\binom{m}{2}\binom{n}{2}$ adet dikdörtgen bulunur.
\end{proof}

\begin{example}[Standart İskambil Eli Sayımı]
$52$ kartlık standart bir iskambil destesinde her biri $13$'er karttan oluşan $4$ simge (Kupa, Karo, Maça, Sinek) ve her simgede $13$ farklı değer ($A, 2, 3, \dots, 10, J, Q, K$) bulunur. Bu desteden seçilen $5$ kartlık bir elde \textbf{tam olarak bir per} (aynı değerden $2$ kart ve diğer $3$ kartın değerleri hem birbirinden hem de perden farklı) bulunma sayısını hesaplayınız.
\end{example}
\begin{proof}[Çözüm]
Sayım adımlarını sırayla oluşturalım:
\begin{enumerate}
  \item \textbf{1. Adım (Perin değerini seçmek):} Perin hangi değerden oluşacağı belirlenir. $13$ farklı değerden biri seçilir: $\binom{13}{1} = 13$ seçenek.
  \item \textbf{2. Adım (Perin kartlarını seçmek):} Seçilen değerin $4$ simgesinden $2$ tanesi belirlenir: $\binom{4}{2} = 6$ seçenek.
  \item \textbf{3. Adım (Kalan 3 kartın değerlerini seçmek):} Elin başka bir per içermemesi için diğer $3$ kartın değerleri birbirinden ve ilk değerden farklı olmalıdır. Kalan $12$ değer arasından $3$ farklı değer seçilir: $\binom{12}{3} = \frac{12 \times 11 \times 10}{6} = 220$ seçenek.
  \item \textbf{4. Adım (Kalan 3 kartın simgelerini seçmek):} Belirlenen $3$ değerin her biri için $4$ simgeden biri seçilir: $\binom{4}{1} \times \binom{4}{1} \times \binom{4}{1} = 4^3 = 64$ seçenek.
\end{enumerate}
Çarpma kuralı gereğince tam olarak bir per içeren ellerin sayısı:
\[ \binom{13}{1} \times \binom{4}{2} \times \binom{12}{3} \times 4^3 = 13 \times 6 \times 220 \times 64 = 1\,098\,240 \]
olarak hesaplanır.
\end{proof}
